\section{Introduction}

% Discuss rise of LLMs in automated system config/API synthesis
% Shift from HitL to agentic workflows/AI agents (ie; x402)
\subsection{Motivation}
In recent years, the deployment of Large Language Models (LLMs) has transitioned
from simple, explanatory or otherwise "conversational" interfaces, to highly
mission-critical integrations with autonomous AI and agentic workflows~\cite{autogpt2023}.
This shift towards "AI agents" and automated system configuration---exemplified
by frameworks such as AutoGPT~\cite{autogpt2023}, OpenClaw~\cite{openclaw2025}, and the x402~\cite{x4022025} protocol---is indicative
of an increasing reliance on these models to synthesize complex API calls and
cloud infrastructure definitions with minimal human-in-the-loop verification
~\cite{openclaw2025}. While the transformer architecture~\cite{vaswani2017attention} has demonstrated
remarkable fluency in natural language, its application in safety-critical
structural synthesis remains somewhat dubious. In industry, the move towards
fully autonomous agents requires that the output not only "looks" like valid
code, but is also actually \textit{compliant} with the operational state of the
target system---that is to say, beyond just syntax~\cite{poesia2022synchromesh}.

% Boundary between CFG/CSI (context-sensitive invariants)
% Why does FSM-based decoding fail? State dependent logic needed?
\subsection{Problem Statement}
Current methodologies for steering LLM output primarily rely on constrained
decoding via Finite State Machines (FSMs) or Context-Free Grammars
(CFGs)~\cite{poesia2022synchromesh}. While these techniques effectively \textit{guarantee}
surface-level syntax (e.g., ensuring a bracket follows a key in a JSON object),
they are fundamentally incapable of enforcing semantic correctness, and any
errors thereof can only be corrected post-hoc, which is computationally costly~\cite{poesia2022synchromesh}.

In other words, there is an important distinction between pure syntactic validity and
semantic correctness that modern approaches do not capture. For instance, an FSM
can ensure a configuration file is "well-formed" JSON, but it cannot verify that
an assigned IP address is within a previously declared subnet range, or that a
\texttt{start\_date} chronologically precedes an \texttt{end\_date}. These are
state-dependent constraints and are therefore beyond the reach of context-free parsers
utilized by modern tools such as Outlines~\cite{willard2023efficient} or Guidance~\cite{lundberg_guidance_2025}. As a consequence,
when LLMs encounter these "semantic walls", they either hallucinate invalid
state transitions, leading to system-wide failures during execution, or post-hoc
checks that demand another prompt, significantly increasing system latency and
cost.

% Introduce dual-process architecture
\subsection{Proposed Contributions}
This work proposes a dual-process neuro-symbolic framework in which a schema-specific 
abstract machine is embedded directly into the LLM's decoding loop to enforce structural
 and semantic constraints, while Oracle-generated traces are used to fine-tune the model 
 and reduce its reliance on runtime intervention.